
\chapter{Conclusion and Recommendations}

In this study, an antireflective coating made of multiple layers of a-Si/SiO$_2$ with a minimum close to the wavelengths at which terahertz photoconductive antennas are optically pumped has been presented. The transfer matrix model was used to obtain design parameters which supplied information for the thicknesses to be used in the fabrication with electron beam physical vapor deposition. The refractive index data used for the transfer matrix model was obtained from measurements on single films.

The thicknesses for a-Si and SiO$_2$ obtained using x-ray reflectivity showed that the quartz crystal microbalance readings set at the reference material parameters generally underestimate the actual thicknesses of the deposited films, but follows a linear relationship with them. The densities of the evaporated films are generally lower than that of the bulk. The refractive index of the a-Si film measured via the envelope method compared to the literature shows a deviation from the usual monotonically decreasing value with increasing wavelength. For the refractive index of the SiO$_2$ film obtained via reflection spectroscopy, the values are higher than the reference, which may be explained by a change in stoichiometric ratio of the film. The measured reflectance of the multilayer coatings generally match the expected minima from the transfer matrix model simulations, although the inaccuracy increases at longer wavelengths.

Recommendations for further studies into the topic include more accurate measurements of the material refractive indices in order to obtain better estimations from the transfer matrix method. Improvements to the quality of the films can also be studied, particularly if substrate heating during growth causes the index to match literature values better. A broader measurement range that includes the infrared region at \qty{1550}{\nano\meter} would also yield more information about the properties of the antireflective coating.

